\documentclass{ctexart}
\usepackage{template}
\usepackage{esint}
\begin{document}\pagestyle{empty}
\begin{center}\Large
北京大学2025-26学年第1学期高数B期末试题
\end{center}

\begin{enumerate}[leftmargin=*,label=\textbf{\arabic*.},topsep=0pt,parsep=0pt,itemsep=0pt,partopsep=0pt]
    \item 试求如下极限.
\begin{enumerate}[label=\tbf{(\alph*)},leftmargin=*,topsep=0pt,parsep=0pt,itemsep=0pt,partopsep=0pt]
\item \[\lim_{x\to0}\dfrac{1-x^2-\e^{-x^2}}{x^2\sin^2x}\]
\item \[\lim_{(x,y)\to(0,0)}\dfrac{xy(x+y)}{x^2+y^2}\]
\end{enumerate}
    \item 设
\[f(x)=\dfrac{\di^3}{\di x^3}(x^2-1)^3\]
\begin{enumerate}[label=\tbf{(\alph*)},leftmargin=*,topsep=0pt,parsep=0pt,itemsep=0pt,partopsep=0pt]
\item 求$f(x)$的显式表达式.
\item 求$f(x)$的单调区间.
\item 求$f(x)$的极值点.
\item 求$f(x)$的凹凸区间.
\end{enumerate}
    \item 已知直线$L_1,L_2$分别为
\[L_1:x=\dfrac y2=\dfrac z3,\qquad L_2:x=-y=z-1\]
若$L_1\subset\pi$, $L_2\parallel\pi$,求平面$\pi$的方程.
    \item 设函数
\[f(x,y)=\begin{cases}1,&x^2+y^2\leqslant1,\\x^2+y^2,&x^2+y^2>1,\end{cases}\qquad \alpha\in(0,\pi/2)\]
证明:在点$(\cos\alpha,\sin\alpha)$处,偏导数$f_x(\cos\alpha,\sin\alpha)$与$f_y(\cos\alpha,\sin\alpha)$均不存在.
    \item 考虑方程
\[F(x,y,z)=xy+yz+\e^{2xz}-4=0\]
\begin{enumerate}[label=\tbf{(\alph*)},leftmargin=*,topsep=0pt,parsep=0pt,itemsep=0pt,partopsep=0pt]
\item 求一点$M_0(x_0,y_0,z_0)$,使得在$M_0$的某邻域内该方程可唯一确定隐函数$z=z(x,y)$,并验证隐函数存在定理的条件.
\item 设在某解点邻域内由方程确定$z=z(x,y)$,求$\dfrac{\p z}{\p x}$, $\dfrac{\p z}{\p y}$.
\end{enumerate}
    \item 设$z=f(u(x,y),v(x,y))$,其中$u(x,y)=\e^x\cos y$, $v(x,y)=\e^x\sin y$.求$\dfrac{\p^2z}{\p y^2}$.
    \item 设$F(x,y)=f(x^2+y^2)$,写出$F(x,y)$在$(x_0,y_0)$附近的二阶Taylor展开式,并用佩亚诺余项表示.
    \item 设$f(x,y,z)=x^2-xy+y^2+z^2$,约束条件为平面$x+y+z=1$.求$f$在该平面上的最小值及其取值点,并证明所得结论.
\end{enumerate}
\end{document}
